\SourceLeaf{101}%
adeoque $qr\equiv\pm f\ (\mathrm{mod.}\ p)$. Combinando hanc congruentiam cum theoremate modo invento obtinemus $rr\equiv\pm2af$, et proin, per artt. 19, 20
\[
2b\equiv\pm rr\ (\mathrm{mod.}\ p)\text{*)}
\]
\begingroup\renewcommand{\thefootnote}{*)}\footnotetext{atque $\{(a\mp b)q\}^2\equiv a\equiv\left(\frac{r-qrr}{2}\right)^2$}\endgroup

Valde memorabile est, discerptionem numeri $p$ in duo quadrata per operationes prorsus directas inveniri posse; scilicet radix quadrati imparis erit residuum absolute minimum ipsius $\frac r{2q}$, radix quadrati paris vero residuum absolute minimum ipsius $\frac12rr$ secundum modulum $p$. Expressionem $\frac r{2q}$, cuius valor pro $p=5$ fit $=1$, pro valoribus maioribus ipsius $p$, ita quoque exhibere licet:
\[
\frac{6\pdot10\pdot14\pdot18\dotsfive(p-3)}{2\pdot3\pdot4\pdot5\dotsfive\frac14(p-1)}
\]

Sed quum insuper noverimus, quonam signo affecta prodeat ex hac formula radix quadrati imparis, eo scilicet, ut semper fiat formae $4m+1$, attentione perdignum est, quod simile criterium generale respectu signi radicis quadrati paris hactenus inveniri non potuerit. Quale si quis inveniat, et nobiscum communicet, magnam de nobis gratiam feret. Interim hic adiungere visum est valores numerorum $a,b,f$, quales pro valoribus ipsius $p$ infra $200$ e residuis minimis expressionum $\frac r{2q}$, $\frac12rr$, $qr$ prodeunt.

\SourceLeaf{102}%
\[
\begin{array}{r|r|r|r}
p&a&b&f\\\hline
5&+1&+2&2\\
13&-3&-2&5\\
17&+1&-4&13\\
29&+5&+2&12\\
37&+1&-6&31\\
41&+5&+4&9\\
53&-7&-2&23\\
61&+5&-6&11\\
73&-3&-8&27\\
89&+5&-8&34\\
97&+9&+4&22\\
101&+1&-10&91\\
109&-3&+10&33\\
113&-7&+8&15\\
137&-11&+4&37\\
149&-7&-10&44\\
157&-11&-6&129\\
173&+13&+2&80\\
181&+9&+10&162\\
193&-7&+12&81\\
197&+1&-14&183
\end{array}
\]
\BookRule{.28\linewidth}

\clearpage
\SourceLeaf{103}%
\thispagestyle{empty}
\pdfbookmark[0]{Theoria residuorum biquadraticorum, commentatio secunda}{paper-five-title}
\begin{center}
\vspace*{.55in}{\large THEORIA\par}
\bigskip{\Large RESIDUORUM BIQUADRATICORUM\par}
\vspace{.4in}{\large COMMENTATIO SECUNDA\par}
\vspace{.6in}AUCTORE\par
\bigskip{\large CAROLO FRIDERICO GAUSS\par}
\vspace{.35in}{\small SOCIETATI REGIAE TRADITA 1831. APR. 15.\par}
\vspace{.55in}\rule{.9\linewidth}{.3pt}\par
\smallskip{\small Commentationes societatis regiae scientiarum Gottingensis recentiores. Vol. \textsc{vii}.\par}
{\small Gottingae \textsc{mdcccxxxii}.\par}
\smallskip\rule{.9\linewidth}{.3pt}
\end{center}

\clearpage
\SourceLeaf{105}%

\renewcommand{\BookRunningTitle}{Theoria residuorum biquadraticorum}
\renewcommand{\Article}[1]{\par\medskip\Needspace{4\baselineskip}\pdfbookmark[1]{Art. #1}{paper-five-art-#1}\centerline{\textbf{#1.}}\nobreak\smallskip}
\begin{center}
\textbf{THEORIA RESIDUORUM BIQUADRATICORUM.}\par
\bigskip COMMENTATIO SECUNDA.
\end{center}
\BookRule{.16\linewidth}
\Article{24}
In commentatione prima ea, quae ad classificationem biquadraticam numeri $+2$ requiruntur, complete absoluta sunt. Dum scilicet omnes numeros per modulum $p$ (qui supponitur esse numerus primus formae $4n+1$) non divisibiles inter quatuor complexus $A,B,C,D$ distributos concipimus, prout singuli ad potestatem exponentis $\frac14(p-1)$ evecti congrui fiunt secundum modulum $p$ ipsi $+1$, $+f$, $-1$, $-f$, denotante $f$ radicem alterutram congruentiae $ff\equiv-1\ (\mathrm{mod.}\ p)$: invenimus, diiudicationem, cuinam complexui adnumerandus sit numerus $+2$, pendere a discerptione numeri $p$ in duo quadrata, ita quidem, ut si statuatur $p=aa+bb$, denotante $aa$ quadratum impar, $bb$ quadratum par, si porro \emph{signa} ipsorum $a,b$ ita accepta supponantur, ut habeatur $a\equiv1\ (\mathrm{mod.}\ 4)$, $b\equiv af\ (\mathrm{mod.}\ p)$, numerus $+2$ ad complexum $A,B,C,D$ pertinere debeat, prout $\frac12b$ sit formae $4n$, $4n+1$, $4n+2$, $4n+3$ resp.

Sponte quoque hinc demanat regula classificationi numeri $-2$ inserviens. Scilicet quum $-1$ pertineat ad classem $A$ pro valore pari ipsius $\frac12b$, ad classem $C$ vero pro impari: pertinebit, per theorema art. 7, numerus $-2$ ad classem $A,B,C,D$, prout $\frac12b$ est formae $4n$, $4n+3$, $4n+2$, $4n+1$ resp.

Haec theoremata etiam sequenti modo exprimi possunt:

\SourceLeaf{106}%
\begin{center}
\begin{tabular}{c|c|c}
Pertinet&$+2$&$-2$\\\hline
ad complexum&\multicolumn{2}{c}{si $b$, secundum modulum $8$, fit congruus ipsi}\\\hline
$A$&$0$&$0$\\
$B$&$2a$&$6a$\\
$C$&$4a$&$4a$\\
$D$&$6a$&$2a$
\end{tabular}
\end{center}

Facile intelligitur, theoremata sic enunciata haud amplius pendere a conditione $a\equiv1\ (\mathrm{mod.}\ 4)$, sed etiamnum valere, si fuerit $a\equiv3\ (\mathrm{mod.}\ 4)$, dummodo conditio altera, $af\equiv b\ (\mathrm{mod.}\ p)$, conservetur.

Aeque facile perspicitur, summam horum theorematum eleganter contrahi posse in formulam unicam, puta:
\begin{quote}\itshape si $a$ et $b$ positive accipiuntur, semper fit\end{quote}
\[
b^{\frac12ab}\equiv a^{\frac12ab}\,2^{\frac14(p-1)}\ (\mathrm{mod.}\ p)
\]

\Article{25}
Videamus nunc, quatenus inductio classificationem numeri $3$ indigitet. Tabula art. 11 ulterius continuata (semper adoptata radice primitiva minima) monstrat, $+3$ pertinere
\begin{center}ad complexum\end{center}
\[
\begin{array}{r|r|r||r|r|r||r|r|r||r|r|r}
\multicolumn{3}{c||}{A\ \text{pro}}&\multicolumn{3}{c||}{B\ \text{pro}}&\multicolumn{3}{c||}{C\ \text{pro}}&\multicolumn{3}{c}{D\ \text{pro}}\\
p&a&b&p&a&b&p&a&b&p&a&b\\
13&-3&+2&17&+1&-4&37&+1&-6&5&+1&+2\\
109&-3&+10&29&+5&+2&61&+5&-6&41&+5&-4\\
181&+9&+10&53&-7&+2&73&-3&-8&149&-7&+10\\
193&-7&-12&89&+5&-8&97&+9&+4&173&+13&+2\\
229&-15&+2&101&+1&+10&157&-11&-6&&&\\
277&+9&+14&113&-7&-8&241&-15&-4&&&\\
&&&137&-11&-4&&&&&&\\
&&&197&+1&-14&&&&&&\\
&&&233&+13&+8&&&&&&\\
&&&257&+1&-16&&&&&&\\
&&&269&+13&+10&&&&&&\\
&&&281&+5&+16&&&&&&\\
&&&293&+17&+2&&&&&&
\end{array}
\]

\SourceLeaf{107}%
Primo saltem aspectu nexum simplicem inter valores numerorum $a,b$, quibus idem complexus respondet, non animadvertimus. At si perpendimus, diiudicationem similem in theoria residuorum quadraticorum per regulam simpliciorem absolvi respectu numeri $-3$, quam respectu numeri $+3$, spes affulget successus aeque secundi in theoria residuorum biquadraticorum. Invenimus autem, $-3$ pertinere ad complexum
\[
\begin{array}{r|r|r||r|r|r||r|r|r||r|r|r}
\multicolumn{3}{c||}{A\ \text{pro}}&\multicolumn{3}{c||}{B\ \text{pro}}&\multicolumn{3}{c||}{C\ \text{pro}}&\multicolumn{3}{c}{D\ \text{pro}}\\
p&a&b&p&a&b&p&a&b&p&a&b\\
37&+1&-6&5&+1&+2&13&-3&+2&29&+5&+2\\
61&+5&-6&17&+1&-4&73&-3&-8&41&+5&-4\\
157&-11&-6&89&+5&-8&97&+9&+4&53&-7&+2\\
193&-7&-12&113&-7&-8&109&-3&+10&101&+1&+10\\
&&&137&-11&-4&181&+9&+10&197&+1&-14\\
&&&149&-7&+10&229&-15&+2&269&+13&+10\\
&&&173&+13&+2&241&-15&-4&293&+17&+2\\
&&&233&+13&+8&277&+9&+14&&&\\
&&&257&+1&-16&&&&&&\\
&&&281&+5&+16&&&&&&
\end{array}
\]
ubi lex inductionis sponte se offert. Scilicet pertinet $-3$ ad complexum
\[
\begin{array}{ll}
A,&\text{quoties }b\text{ per }3\text{ divisibilis est, sive }b\equiv0\ (\mathrm{mod.}\ 3)\\
B,&\text{quoties }a+b\text{ per }3\text{ est divisibilis, sive }b\equiv2a\ (\mathrm{mod.}\ 3)\\
C,&\text{quoties }a\text{ per }3\text{ est divisibilis, sive }a\equiv0\ (\mathrm{mod.}\ 3)\\
D,&\text{quoties }a-b\text{ per }3\text{ est divisibilis, sive }b\equiv a\ (\mathrm{mod.}\ 3)
\end{array}
\]

\Article{26}
Numerum $+5$ adscribendum invenimus complexui
\[
\begin{array}{ll}
A\text{ pro }p=&101,\ 109,\ 149,\ 181,\ 269\\
B\text{ pro }p=&13,\ 17,\ 73,\ 97,\ 157,\ 193,\ 197,\ 233,\ 277,\ 293\\
C\text{ pro }p=&29,\ 41,\ 61,\ 89,\ 229,\ 241,\ 281\\
D\text{ pro }p=&37,\ 53,\ 113,\ 137,\ 173,\ 257
\end{array}
\]
In considerationem vocatis valoribus numerorum $a,b$ singulis $p$ respondentibus, lex hic aeque facile, ut pro classificatione numeri $-3$, prehenditur. Scilicet incidimus in complexum

\SourceLeaf{108}%
\[
\begin{array}{ll}
A,&\text{quoties }b\equiv0\ (\mathrm{mod.}\ 5)\\
B,&\text{quoties }b\equiv a\\
C,&\text{quoties }a\equiv0\\
D,&\text{quoties }b\equiv4a
\end{array}
\]

Manifestum est, has regulas complecti casus omnes, quum pro $b\equiv2a$, vel $b\equiv3a\ (\mathrm{mod.}\ 5)$, fieret $aa+bb\equiv0$, Q. E. A., quum per hypothesin $p$ sit numerus primus a $5$ diversus.

\Article{27}
Perinde inductio ad numeros $-7$, $-11$, $+13$, $+17$, $-19$, $-23$ applicata satisque producta sequentes regulas indigitat:
\begin{center}Pro numero $-7$\end{center}
\[
\begin{array}{c|l}
A&a\equiv0,\quad\text{vel }b\equiv0\ (\mathrm{mod.}\ 7)\\
B&b\equiv4a,\ \text{vel }b\equiv5a\\
C&b\equiv a,\quad\text{vel }b\equiv6a\\
D&b\equiv2a,\ \text{vel }b\equiv3a
\end{array}
\]

\begin{center}Pro numero $-11$.\end{center}
\[
\begin{array}{c|l}
A&b\equiv0,\ 5a,\ \text{vel }6a\ (\mathrm{mod.}\ 11)\\
B&b\equiv a,\ 3a\text{ vel }4a\\
C&a\equiv0,\ \text{vel }b\equiv2a\text{ vel }9a\\
D&b\equiv7a,\ 8a\text{ vel }10a
\end{array}
\]

\begin{center}Pro numero $+13$.\end{center}
\[
\begin{array}{c|l}
A&b\equiv0,\ 4a,\ 9a\ (\mathrm{mod.}\ 13)\\
B&b\equiv6a,\ 11a,\ 12a\\
C&a\equiv0;\ b\equiv3a,\ 10a\\
D&b\equiv a,\ 2a,\ 7a
\end{array}
\]

\begin{center}Pro numero $+17$.\end{center}
\[
\begin{array}{c|l}
A&a\equiv0;\ b\equiv0,\ a,\ 16a\ (\mathrm{mod.}\ 17)\\
B&b\equiv2a,\ 6a,\ 8a,\ 14a\\
C&b\equiv5a,\ 7a,\ 10a,\ 12a\\
D&b\equiv3a,\ 9a,\ 11a,\ 15a
\end{array}
\]

\SourceLeaf{109}%
\begin{center}Pro numero $-19$.\end{center}
\[
\begin{array}{c|l}
A&b\equiv0,\ 2a,\ 5a,\ 14a,\ 17a\ (\mathrm{mod.}\ 19)\\
B&b\equiv3a,\ 7a,\ 11a,\ 13a,\ 18a\\
C&a\equiv0;\ b\equiv4a,\ 9a,\ 10a,\ 15a\\
D&b\equiv a,\ 6a,\ 8a,\ 12a,\ 16a
\end{array}
\]

\begin{center}Pro numero $-23$.\end{center}
\[
\begin{array}{c|l}
A&a\equiv0;\ b\equiv0,\ 7a,\ 10a,\ 13a,\ 16a\ (\mathrm{mod.}\ 23)\\
B&b\equiv2a,\ 3a,\ 4a,\ 11a,\ 15a,\ 17a\\
C&b\equiv a,\ 5a,\ 9a,\ 14a,\ 18a,\ 22a\\
D&b\equiv6a,\ 8a,\ 12a,\ 19a,\ 20a,\ 21a
\end{array}
\]

\Article{28}
Theoremata specialia hoc modo per inductionem eruta confirmari inveniuntur, quousque haec continuetur, formamque criteriorum pulcherrimam manifestant. Si vero inter se conferuntur, ut conclusiones generales inde petantur, primo statim aspectu se offerunt observationes sequentes.

Criteria diiudicationis, ad quemnam complexum referendus sit numerus primus $\pm q$ (sumendo signum superius vel inferius, prout $q$ est formae $4n+1$ vel $4n+3$), pendent a formis numerorum $a,b$ inter se collatorum respectu moduli $q$. Scilicet

I. quoties $a\equiv0\ (\mathrm{mod.}\ q)$, $\pm q$ pertinet ad complexum determinatum, qui est $A$ pro $q=7,\ 17,\ 23$, nec non $C$ pro $q=3,\ 11,\ 13,\ 19$, unde coniectura oritur, casum priorem generaliter valere, quoties $q$ sit formae $8n\pm1$, posteriorem vero, quoties $q$ sit formae $8n\pm3$. Ceterum complexus $B$ et $D$ iam absque inductione excluduntur pro valore ipsius $a$ per $q$ divisibili, ubi fit $p\equiv bb\ (\mathrm{mod.}\ q)$, i. e. ubi $p$ est residuum quadraticum ipsius $q$, unde per theorema fundamentale $\pm q$ esse debet residuum quadraticum ipsius $p$.

II. Quoties autem $a$ per $q$ non est divisibilis, criterium pendet a valore expressionis $\frac ba\ (\mathrm{mod.}\ q)$. Admittit quidem haec expressio $q$ valores diversos, puta $0,\ 1,\ 2,\ 3\dotsfour q-1$: sed quoties $q$ est formae $4n+1$, excludendi sunt bini valo%
\SourceLeaf{110}%
res expressionis $\surd-1\ (\mathrm{mod.}\ q)$, qui manifesto nequeunt esse valores expressionis $\frac ba\ (\mathrm{mod.}\ q)$, quum $p=aa+bb$ semper supponatur esse numerus primus a $q$ diversus. Quapropter multitudo valorum admissibilium expressionis $\frac ba\ (\mathrm{mod.}\ q)$ est $=q-2$, pro $q\equiv1\ (\mathrm{mod.}\ 4)$, dum manet $=q$ pro $q\equiv3\ (\mathrm{mod.}\ 4)$.

Iam hi valores in quaternas classes distribuuntur, puta, ut quidam, indefinite per $\alpha$ denotandi, respondeant complexui $A$; alii per $\beta$ denotandi complexui $B$; alii $\gamma$ complexui $C$; denique reliqui $\delta$ complexui $D$, ita scilicet, ut $\pm q$ complexui $A,B,C,D$ adscribendus sit, prout habeatur $b\equiv\alpha a$, $b\equiv\beta a$, $b\equiv\gamma a$, $b\equiv\delta a\ (\mathrm{mod.}\ q)$.

At \emph{lex} huius distributionis abstrusior videtur, etiamsi quaedam generalia promte animadvertantur. Multitudo in ternis classibus eadem reperitur, puta $=\frac14(q-1)$ vel $\frac14(q+1)$, dum in una (et quidem in eadem, quae respondet complexui cum criterio $a\equiv0$) unitate minor est, ita ut multitudo omnium criteriorum diversorum respectu singulorum complexuum fiat eadem, puta $=\frac14(q-1)$ vel $\frac14(q+1)$. Porro animadvertimus, $0$ semper in prima classe (inter $\alpha$) reperiri, nec non complementa numerorum $\alpha,\beta,\gamma,\delta$ ad $q$, puta $q-\alpha$, $q-\beta$, $q-\gamma$, $q-\delta$ resp. in classe prima, quarta, tertia, secunda. Denique valores expressionum $\frac1\alpha$, $\frac1\beta$, $\frac1\gamma$, $\frac1\delta\ (\mathrm{mod.}\ q)$ pertinere videmus ad classem primam, quartam, tertiam, secundam, quoties criterium $a\equiv0$ respondet complexui $A$; ad classem tertiam, secundam, primam, quartam resp. autem, quoties criterium $a\equiv0$ refertur ad complexum $C$. Sed ad haec fere limitantur, quae per inductionem assequi licet, nisi audacius ea, quae infra e fontibus genuinis haurientur, anticipare nobis arrogemus.

\Article{29}
Antequam ulterius progrediamur, observare convenit, criteria pro numeris primis (positive sumtis, si sunt formae $4n+1$, negative, si formae $4n+3$) sufficere ad diiudicationem pro omnibus reliquis numeris, si modo theorema art. 7, atque criteria pro $-1$ et $\pm2$ in subsidium vocentur. Ita e. g. si desiderantur criteria pro numero $+3$, criteria in art. 25 prolata, quae referuntur ad $-3$, etiamnum pro $+3$ valebunt, quoties $\frac12b$ est numerus par: contra complexus $A,B,C,D$ cum complexibus $C,D,A,B$ permutandi erunt, quoties $\frac12b$ est impar, unde sequuntur praecepta haecce:

